% B3PoW-Scratch v1.1 -- IACR ePrint preprint.
%
% Target style: Springer Lecture Notes in Computer Science (`llncs.cls`).
% This file does NOT use llncs by default because `llncs.cls` is not part
% of a standard TeX Live install -- it ships as a separate download from
% Springer (https://www.springer.com/gp/computer-science/lncs/conference-
% proceedings-guidelines). To keep the build reproducible on a vanilla
% TeX Live 2023+, we fall back to the `article` class and re-implement
% the small subset of LNCS conveniences we use (\institute, \keywords).
%
% To target LNCS for the actual ePrint submission:
%
%   1. Drop `llncs.cls` into this directory (or install it system-wide).
%   2. Replace the `\documentclass{...}` line below with
%        \documentclass[runningheads]{llncs}
%   3. Drop the small "LNCS-shim" block (search for: LNCS-SHIM-BEGIN).
%   4. Change the \bibliographystyle to `splncs04` (already commented in).
%
% Build:
%   make           # -> B3POW-SCRATCH-PREPRINT.pdf  (latexmk -xelatex -bibtex)
%   make clean
%   make check     # lacheck + chktex, if installed.

\documentclass[11pt,a4paper]{article}

% --- Packages (mirror the LNCS defaults where they overlap) ---------------
% Targets XeLaTeX (Unicode-native); inputenc/fontenc deliberately omitted.
\usepackage[a4paper,margin=1in]{geometry}
\usepackage{amsmath,amsthm,amssymb}
\usepackage{graphicx}
\usepackage{url}
\usepackage{booktabs}
\usepackage{listings}
\usepackage{microtype}
\usepackage{xcolor}
\usepackage[hidelinks,breaklinks=true]{hyperref}
\usepackage{algorithm}
\usepackage{algpseudocode}
\usepackage{caption}

% --- LNCS-SHIM-BEGIN ------------------------------------------------------
% Re-implements just enough of llncs.cls's interface that the source below
% reads as if it were LNCS. Delete this block when switching to llncs.
% \institute{...} is a no-op in article mode; the affiliation is embedded
% directly in \author{...} below. llncs.cls provides the real \institute.
\providecommand{\institute}[1]{}
\providecommand{\keywords}[1]{\par\medskip\noindent\textbf{Keywords:} #1.}
% Theorem environments -- LNCS provides these unnumbered-by-default; we
% number them within section for paper-readability.
\theoremstyle{plain}
\newtheorem{theorem}{Theorem}[section]
\newtheorem{lemma}[theorem]{Lemma}
\newtheorem{proposition}[theorem]{Proposition}
\theoremstyle{definition}
\newtheorem{definition}[theorem]{Definition}
\theoremstyle{remark}
\newtheorem{remark}[theorem]{Remark}
% --- LNCS-SHIM-END --------------------------------------------------------

% --- Listings: a Python-flavour skin used very sparingly ------------------
\definecolor{lstbg}{gray}{0.97}
\lstset{
  basicstyle=\footnotesize\ttfamily,
  backgroundcolor=\color{lstbg},
  breaklines=true,
  showstringspaces=false,
  frame=single,
  framerule=0.4pt,
  framesep=4pt,
  xleftmargin=0pt,
  xrightmargin=0pt,
  language=Python,
  keywordstyle=\bfseries,
  commentstyle=\itshape\color{black!55},
  stringstyle=\color{black!70}
}

% --- Title block ----------------------------------------------------------
% Under article mode (the default here) the affiliation is folded into
% \author{}. Under llncs.cls the \institute{...} line below the document
% start takes over; the no-op shim above keeps that line harmless in
% article mode.
\title{B3PoW-Scratch v1.1:\\
       A Memory-Hard BLAKE3-Based Proof-of-Work\\
       with On-Chip-Memory-Bounded Hardware Economy}
\author{The B3Chain Project\thanks{Correspondence:
        \href{mailto:paper@b3chain.org}{\texttt{paper@b3chain.org}}.
        Authorship placeholder; final list to be confirmed at
        ePrint submission.} \\[2pt]
        \small Open-source project, \url{https://b3chain.org} \\
        \small Repository: \url{https://github.com/b3chain/b3chain}}
\date{}

\begin{document}

\maketitle

% The next line is a no-op under article mode (see the shim) and the
% real LNCS affiliation block under llncs.cls.
\institute{Open-source project, \url{https://b3chain.org}.
           Source repository at \url{https://github.com/b3chain/b3chain}.}

% --- Abstract -------------------------------------------------------------
\begin{abstract}
We describe B3PoW-Scratch~v1.1, the memory-hard Proof-of-Work used by
B3Chain, a conservative Bitcoin Core fork that keeps Bitcoin's UTXO
model, monetary schedule, block-identity hash, and P2P protocol and
replaces only the mining hash function. B3PoW-Scratch is built on
BLAKE3 over a 1\,MiB scratchpad partitioned into eight parallel lanes
and executes 2\,048 sequential data-dependent read--modify--write
iterations per hash. The construction is calibrated, not theatrical:
the design target is to make the most economical production miner a
small on-chip-memory FPGA card rather than a custom ASIC or a
commodity GPU. We make no claim of permanent ASIC resistance. We
evaluate the construction against the catalogued PoW attack
surface---preimage and collision, memory-hardness shortcuts, address
uniformity, verification-DoS, 51\,\%~chain reorganisation, and
time-warp---and against four hardware tiers (CPU, GPU, FPGA, ASIC).
A four-implementation parity gate (Python, C++, TypeScript,
SystemVerilog) enforces byte-equivalence across an in-tree consensus
vector set, removing the most common silent-failure mode for new PoW
launches.
\end{abstract}

\keywords{Proof of work \and memory-hard hashing \and BLAKE3
          \and FPGA mining \and blockchain consensus}

% --- 1. Introduction ------------------------------------------------------
\section{Introduction}\label{sec:intro}

A new Bitcoin-derived Proof-of-Work (PoW) blockchain that adopts
Bitcoin's SHA-256d mining hash inherits, by construction, the global
SHA-256d ASIC fleet on day one. At any plausible launch hashrate the
idle capacity of that fleet exceeds the new chain's honest hashrate
by many orders of magnitude, and a single pool operator can rewrite
the chain at a marginal electricity cost. The SHA-256d
\emph{carry-over problem} is the dominant failure mode for new
Bitcoin-derived chains in their first months. Choosing a memory-hard
or GPU-friendly algorithm avoids it but trades it for different
failure modes: CPU/GPU-friendly designs invite botnet hashrate,
designs that depend on exotic memory hierarchies are unverifiable on
commodity nodes, and designs that are too hardware-agnostic recreate
the SHA-256d race on a new vertex.

B3PoW-Scratch~v1.1 takes a calibrated position on this trade-off
curve. Its explicit design target is a small, low-power FPGA card
with \emph{on-chip} memory~--- specifically, the Xilinx Kintex
UltraScale+ XCKU5P~\cite{xilinx-ku5p}~--- as the most economical
production miner. The chain is otherwise Bitcoin: 21\,M cap,
10-minute target, 210\,k halving, all Bitcoin script opcodes,
SHA-256d block identity, Bitcoin's wire format. Only the PoW hash
function differs.

\paragraph{Contributions.} This paper:
\begin{itemize}\itemsep0.2em
\item formalises B3PoW-Scratch~v1.1 as a sequential read--modify--write
      memory-hard hash on a 1\,MiB scratchpad with eight parallel
      lanes and 2\,048 iterations
      (\S\ref{sec:algorithm});
\item argues that the construction inherits BLAKE3's collision /
      preimage resistance, exhibits a non-amortisable sequential
      memory-bandwidth lower bound similar to that of
      \texttt{scrypt}~\cite{percival2009scrypt} and
      Argon2-d~\cite{biryukov2016argon2}, and is hardened against the
      catalogued PoW attack surface (\S\ref{sec:security});
\item evaluates hardware economy at four tiers (CPU, GPU, FPGA, ASIC)
      and shows that the algorithm's combination of working-set size
      and sequential dependency chain pins the economical operating
      point to a small on-chip-memory FPGA at the launch hashrate
      regime (\S\ref{sec:hardware});
\item documents a four-implementation in-tree parity gate (Python,
      C++, TypeScript, SystemVerilog) and the consensus-vector
      regime that enforces byte-equivalence across them
      (\S\ref{sec:impl}); and
\item enumerates the open questions we are aware of, including those
      that we explicitly defer to external review
      (\S\ref{sec:openq}).
\end{itemize}

\paragraph{What we do not claim.} We do not claim a novel
cryptographic primitive; BLAKE3 is used unmodified as the only
non-linear mixing primitive. We do not claim permanent
ASIC-resistance; \S\ref{sec:hardware} bounds the per-die advantage of
a custom ASIC over the reference FPGA at single-digit-to-low-double-digit
multipliers, not at \emph{infinity}. We do not claim a finished
security analysis; the arguments here are informal, an external audit
is scoped~\cite{b3chain-auditscope}, and this preprint exists in part
to solicit reviewer attention.

% --- 2. Related work ------------------------------------------------------
\section{Background and Related Work}\label{sec:related}

\paragraph{SHA-256d-based chains.}
Bitcoin's SHA-256d~\cite{nakamoto2008bitcoin} is compute-bound and
arithmetically simple, properties that made it ASIC-able within
roughly three years of mainnet launch. The economic consequence is
the carry-over problem stated above: any new chain that adopts the
same algorithm inherits the existing fleet rather than recruiting
hardware. Algorithm-rotated forks (Litecoin, Bitcoin~Gold, Bitcoin
Cash variants) trade this for new equilibria with their own
asymmetries.

\paragraph{Memory-hard PoW lineage.}
Scrypt~\cite{percival2009scrypt} introduced memory-hardness to
crypto-economic key derivation and was adopted by Litecoin as PoW; in
practice its small parameterised working set proved ASIC-able. The
Argon2 family~\cite{biryukov2016argon2}, RFC~9106~\cite{rfc9106},
introduced data-dependent (Argon2d), data-independent (Argon2i), and
hybrid (Argon2id) variants targeted at password hashing; its
throughput--verification trade-offs are wrong for PoW economics in
isolation (verification cost scales with the memory parameter).
RandomX~\cite{tevador2019randomx} targets the CPU as the
economically optimum execution model via program generation and a
256\,MB dataset; its verification cost is non-trivial and its
ASIC-resistance argument is informal.
ProgPoW~\cite{programmaticpow2019} targets GPU-economical operation
via program generation; it was discussed for Ethereum but never
deployed on the canonical chain. Equihash~\cite{biryukov2017equihash}
applies Wagner's generalised birthday algorithm~\cite{wagner2002gbp}
as a memory-asymmetric PoW; it was ASIC-ed at the lower memory
parameter sets at which Zcash deployed it. Ethash and
KawPow~\cite{ethash} are GPU-targeted PoWs built around a slowly
growing DAG.

\paragraph{BLAKE3 as the primitive.}
BLAKE3~\cite{oconnor2020blake3} is a 256-bit cryptographic hash
function descended from the BLAKE2 family~\cite{aumasson2013blake2},
itself a SHA-3 finalist. Its inner round (the \texttt{g}-function)
is hardware-friendly: a Xilinx 7-series or later instantiation of one
\texttt{g} fits in roughly 600 LUTs, which is what made BLAKE3 a
natural choice for an FPGA-targeted PoW.

\paragraph{Verification-DoS and time-warp prior art.}
Two known classes of attack on the validation side inform our design:
verification-DoS, where an adversary submits headers crafted to be
expensive to verify (the broader category that includes the
fix-it-fast incident underlying CVE-2018-17144~\cite{cve2018-17144}
in a related context); and time-warp, where an adversary manipulates
header timestamps to depress difficulty within a retarget window
(closed on Bitcoin mainnet by BIP-94~\cite{bip94}). Our verifier
mitigations (\S\ref{sec:security}) borrow directly from these
prior fixes. Selfish mining~\cite{eyal2014selfish} (and its
broader analyses~\cite{bonneau2015sok,sapirshtein2016selfish})
defines the lower attack-share threshold below which the canonical
51\,\% formula understates risk.

\paragraph{Position.} The axis along which B3PoW-Scratch differs
from the prior memory-hard PoWs is the \emph{target execution
model}: RandomX optimises for the CPU; ProgPoW and KawPow for the
GPU; B3PoW-Scratch for a small FPGA whose entire scratchpad fits
on-chip in BRAM. That choice shapes every other parameter
(working-set size, lane structure, verification-cost ceiling). The
construction is deliberately uninteresting cryptographically; its
contribution is the calibration.

% --- 3. Design goals & threat model ---------------------------------------
\section{Design Goals and Threat Model}\label{sec:goals}

\subsection{Design goals}\label{sec:goals-d}

We state the design goals as quantitative targets so they are
falsifiable:

\begin{enumerate}\itemsep0.2em
\item[\textbf{G1}] \emph{On-chip-memory-bound.} The working set fits
      entirely in the on-chip memory (BRAM~+ URAM, no external DRAM)
      of the reference target Xilinx XCKU5P~\cite{xilinx-ku5p}.
      We choose 1\,MiB so that an 8-lane partition occupies
      $\approx 35\%$ of the KU5P's 432 BRAM tiles, leaving room for
      mixer pipelines, FIFOs, and control logic.
\item[\textbf{G2}] \emph{Sequential RMW chain.} Each iteration's
      address derivation, scratchpad read, lane mixing, and
      scratchpad writeback depend on the previous iteration's lane
      state. There is no Bitcoin-style midstate or parallel-nonce
      shortcut inside a single hash.
\item[\textbf{G3}] \emph{BLAKE3 primitive only.} The only non-linear
      mixing primitive used is the standard BLAKE3 round function;
      the construction's security reduces to BLAKE3's published
      claims.
\item[\textbf{G4}] \emph{Verification cost $<1\%$ of block interval.}
      A single verification of one PoW hash on a 2026-era CPU core
      must complete in well under 50\,ms; the 10-minute target
      interval gives us margin to spare.
\item[\textbf{G5}] \emph{Bounded ASIC advantage.} A custom ASIC that
      implements B3PoW-Scratch will outperform the reference FPGA;
      the design intent is to bound that advantage to single-digit /
      low-double-digit multipliers, not the $10^6$ multiplier that
      SHA-256d ASICs enjoy over CPUs. \S\ref{sec:hardware} makes the
      bound quantitative.
\end{enumerate}

\subsection{Threat model}\label{sec:goals-t}

We protect against:

\begin{itemize}\itemsep0.2em
\item \emph{An attacker with idle SHA-256d ASIC capacity.} Trivially
      neutralised by the algorithm change: SHA-256d ASICs cannot
      mine B3PoW-Scratch at any difficulty.
\item \emph{An attacker with significant GPU capacity.}
      \S\ref{sec:hardware} argues that GPU H/s\,/\,\$ on
      B3PoW-Scratch is below CPU H/s\,/\,\$, eliminating most
      rational GPU-attack scenarios.
\item \emph{An attacker who builds a custom B3PoW-Scratch ASIC.}
      \S\ref{sec:hardware} bounds the per-die advantage.
\item \emph{A node-validator-DoS attacker} who crafts headers to
      force expensive validation. Mitigated by a per-hash wall-clock
      budget, a per-parent pad cache, and the headers-first sync
      pattern inherited from Bitcoin Core
      (\S\ref{sec:security-verifydos}).
\item \emph{A 51\,\%-style chain-reorganisation attacker} at all
      hashrate shares. Full quantitative treatment is in the
      project's 51\,\%-attack analysis~\cite{b3chain-51attack};
      \S\ref{sec:security-51} summarises the calibrated bounds.
\end{itemize}

We do \emph{not} protect against:

\begin{itemize}\itemsep0.2em
\item A state-level attacker with effectively unbounded capital
      against a young low-hashrate chain. (True of every PoW chain.)
\item A successful cryptanalytic break of BLAKE3. (Inherited
      bound.)
\item Compromised firmware on the miner hardware, which is
      addressed by signed-manifest verification outside the
      consensus layer.
\end{itemize}

% --- 4. Algorithm ---------------------------------------------------------
\section{Algorithm}\label{sec:algorithm}

This section is informative; the byte-exact normative
reference is the project's \texttt{SPEC.md}~\cite{b3chain-spec}, and
the executable reference is \texttt{ref/b3pow\_ref.py} in the
b3chain repository. The four in-tree implementations
(\S\ref{sec:impl}) are gated against a single consensus-vector set.

\subsection{Parameters}\label{sec:algo-params}

The construction is parameterised by the constants in
Table~\ref{tab:params}, all of which are consensus-critical and
encoded in \texttt{SPEC\_VERSION}~$=$~\texttt{0x00010101}.

\begin{table}[t]
\centering
\caption{B3PoW-Scratch v1.1 parameters (KU5P profile). Source:
         SPEC~\S3~\cite{b3chain-spec}.}
\label{tab:params}
\begin{tabular}{lll}
\toprule
Symbol & Value & Justification \\
\midrule
\texttt{SCRATCH\_BYTES}  & $1\,048\,576$        & 1\,MiB; fits in
                                                 KU5P on-chip BRAM \\
\texttt{LANES}           & $8$                  & one BRAM partition
                                                 per lane \\
\texttt{LANE\_BYTES}     & $131\,072$           & 128\,KiB per lane \\
\texttt{BLOCK\_BYTES}    & $64$                 & BLAKE3 block size \\
\texttt{LANE\_BLOCKS}    & $2\,048$             & 11-bit address window
                                                 per lane \\
\texttt{ITERATIONS}      & $2\,048$             & latency dominates
                                                 init cost \\
\texttt{INNER\_ROUNDS}   & $2$                  & reduced BLAKE3
                                                 rounds per mix \\
\texttt{BLAKE3\_ROUNDS}  & $7$                  & full BLAKE3 for
                                                 init / seed / final \\
\bottomrule
\end{tabular}
\end{table}

The 8 multiplicative-mixer constants
$\texttt{ITER\_MUL}[0..7]\in\mathbb{Z}_{2^{64}}$ are pairwise distinct
wyhash-vetted secrets~\cite{wyhash}; the v1.1.1 release
(\texttt{SPEC\_VERSION = 0x00010101}) restored distinctness after a
duplicate was discovered in v1.1.0 (finding F-1
in~\cite{b3chain-51attack}; see~\S\ref{sec:security-uniformity}).

\subsection{Scratchpad initialisation}\label{sec:algo-init}

For a parent block hash $p\in\{0,1\}^{256}$ the scratchpad
$P\in\{0,1\}^{8\,388\,608}$ is derived block-wise using
BLAKE3-XOF~\cite{oconnor2020blake3} in extendable-output mode:
$$
P[64i, 64(i{+}1)) \;=\; \mathrm{BLAKE3\text{-}XOF}\!\bigl(p \,\Vert\,
\mathrm{LE32}(i)\bigr)\quad \text{for } i\in[0,16\,384).
$$
$P$ depends only on $p$, so it is shared across all nonce attempts
for a given parent. A 1\,MiB pad is roughly 100 BLAKE3 chunk
compressions and dominates the per-parent cost (the per-nonce cost
is the main loop).

\subsection{Lane setup}\label{sec:algo-lanes}

For a serialised header $H\in\{0,1\}^{640}$ (the standard 80-byte
Bitcoin block header), the per-lane initial 256-bit state is:
$$
\textsc{lanes}[L] \;=\; \mathrm{BLAKE3}\!\bigl(\mathrm{BLAKE3}(H)
\,\Vert\, \mathrm{LE32}(L)\bigr) \qquad L\in[0,8).
$$

\subsection{Address derivation}\label{sec:algo-addr}

The per-lane scratchpad block address for iteration $r$ is computed
from the lane state by a multiplicatively-keyed mix:

\begin{algorithm}[t]
\caption{\textsc{DeriveAddresses}; the per-iteration per-lane address
function.}
\label{alg:addr}
\begin{algorithmic}[1]
\Require lanes $\in (\{0,1\}^{256})^8$, iteration index
         $r\in [0, 2\,048)$
\Ensure addresses $\textsc{addr}[L]\in [0, 2\,048)$ for
        $L\in[0,8)$
\For{$L = 0, \ldots, 7$}
  \State $\textsc{lo} \gets $ first 8 bytes of lanes$[L]$,
         little-endian as $\mathbb{Z}_{2^{64}}$
  \State $\textsc{hi} \gets $ next 8 bytes of lanes$[L]$,
         little-endian as $\mathbb{Z}_{2^{64}}$
  \State $\textsc{mul} \gets ((\textsc{hi} \oplus r) \cdot
         \texttt{ITER\_MUL}[L]) \bmod 2^{64}$
  \State $\textsc{mixed} \gets \textsc{lo} \oplus
         \mathrm{rotr}_{64}(\textsc{mul}, 23)$
  \State $\textsc{addr}[L] \gets \textsc{mixed} \bmod 2\,048$
\EndFor
\State \Return $\textsc{addr}$
\end{algorithmic}
\end{algorithm}

The iteration index $r$ is XOR-mixed into \emph{the high half of} the
lane state \emph{before} the multiplication, so that its bits
diffuse through every output bit (multiplication mixes high and low
halves of its operands). The post-rotate alternative leaves $r$ at
bit position~41 of \textsc{mixed} where it never reaches the 11-bit
address window.

\subsection{Mix step}\label{sec:algo-mix}

Each iteration applies a reduced-round BLAKE3 compression to the
(lane-state, scratchpad-block) pair and a lane permutation
$\sigma\colon L\mapsto (5L+1)\bmod 8$ to the resulting lane states.
Let $g$ denote the BLAKE3 quarter-round mixer
(\cite{oconnor2020blake3}, \S2.5) and $\pi$ the standard 16-word
message permutation. The per-lane mix is then:

\begin{algorithm}[t]
\caption{\textsc{MixStep}; one iteration's lane-local mixing.
\texttt{INNER\_ROUNDS}~$=2$.}
\label{alg:mix}
\begin{algorithmic}[1]
\Require lanes $\in (\{0,1\}^{256})^8$,
         blocks $\in (\{0,1\}^{512})^8$
\Ensure $(\text{lanes}', \text{blocks}')$
\For{$L = 0, \ldots, 7$}
  \State $s \gets \text{lanes}[L] \,\Vert\, \texttt{IV}[0..3]
         \,\Vert\, 0, 0, 64, 0$
         \Comment{16 32-bit words}
  \State $m \gets \text{blocks}[L]$ as 16 little-endian 32-bit words
  \For{$j = 1, \ldots, \texttt{INNER\_ROUNDS}$}
    \State apply $g$ to the 4 columns and 4 diagonals of $s$ using
           $m$ as message
    \State $m \gets \pi(m)$
  \EndFor
  \State $\text{lanes}'[L] \gets [s_i \oplus s_{i+8}]_{i=0}^{7}$
  \State $\text{blocks}'[L] \gets \text{blocks}[L] \oplus
         \mathrm{pack}_\text{LE}(m)$
\EndFor
\State $\text{lanes}' \gets [\text{lanes}'[\sigma(L)]]_{L=0}^{7}$
\Comment{lane shuffle}
\State \Return $(\text{lanes}', \text{blocks}')$
\end{algorithmic}
\end{algorithm}

Two rounds is the minimum at which BLAKE3's $g$ yields full
intra-state diffusion on the cv--message pair, and the lane shuffle
$\sigma$ provides inter-lane diffusion within a small number of
iterations (Lemma~\ref{lem:diffusion}).

\subsection{Top-level pipeline}\label{sec:algo-top}

\begin{algorithm}[t]
\caption{\textsc{B3PoW-Scratch}; the top-level hash.}
\label{alg:top}
\begin{algorithmic}[1]
\Require header $H\in\{0,1\}^{640}$, parent hash $p\in\{0,1\}^{256}$
\Ensure 256-bit PoW hash
\State seed $\gets \mathrm{BLAKE3}(H)$
\State $P \gets \textsc{InitScratchpad}(p)$
\Comment{cacheable per-parent}
\State lanes $\gets \textsc{InitLanes}(\text{seed})$
\For{$r = 0, \ldots, 2\,047$}
  \State addr $\gets \textsc{DeriveAddresses}(\text{lanes}, r)$
  \State blocks $\gets \textsc{Read}(P, \text{addr})$
  \State $(\text{lanes}, \text{blocks}') \gets
         \textsc{MixStep}(\text{lanes}, \text{blocks})$
  \State $\textsc{Write}(P, \text{addr}, \text{blocks}')$
\EndFor
\State \Return $\mathrm{BLAKE3}(\text{lanes}[0] \Vert \cdots \Vert
       \text{lanes}[7] \Vert H[76, 80))$
\end{algorithmic}
\end{algorithm}

\subsection{PoW check}\label{sec:algo-check}

The hash is accepted iff it satisfies the unsigned little-endian
256-bit inequality
\begin{equation}\label{eq:powcheck}
\mathrm{int}_{\mathrm{LE}}\bigl(\textsc{B3PoW-Scratch}(H, p)\bigr)
\;\le\; \tau(\texttt{nBits}),
\end{equation}
where $\tau$ is Bitcoin's compact-encoded target unchanged from
\texttt{arith\_uint256::SetCompact}.

\subsection{Per-hash dimensions}\label{sec:algo-dims}

A single hash invokes BLAKE3 in three distinct modes: $16\,384$
BLAKE3-XOF chunk compressions for the scratchpad init, $9$ full
BLAKE3 hashes for the seed and lane setup, $16\,384$ reduced
two-round compressions inside the inner loop, and one final
BLAKE3. Memory traffic per hash is $\approx 8$\,MiB read and
$\approx 8$\,MiB write~--- sequential within a lane, scattered
across lanes by the address derivation.

% --- 5. Security analysis -------------------------------------------------
\section{Security Analysis}\label{sec:security}

The arguments below are informal; the project's external audit
scope~\cite{b3chain-auditscope} requests formal verification of each.

\subsection{Preimage and collision resistance}\label{sec:security-pc}

\begin{theorem}\label{thm:inherited}
If BLAKE3 is collision-resistant with security level $\lambda_c$ and
preimage-resistant with security level $\lambda_p$, then
B3PoW-Scratch~v1.1 is collision-resistant with security $\lambda_c$
and preimage-resistant with security $\lambda_p$.
\end{theorem}

\begin{proof}[Sketch]
The output of \textsc{B3PoW-Scratch} is a single application of
$\mathrm{BLAKE3}$ to a 260-byte input (the eight 32-byte lane states
concatenated with the 4-byte nonce). Any collision (preimage) on
B3PoW-Scratch yields a collision (preimage) on BLAKE3 applied to the
260-byte lanes-and-nonce buffer~--- the inner loop influences only
the value of that buffer, not the BLAKE3 evaluation performed on
it. Standard claims for BLAKE3~\cite{oconnor2020blake3} give
$\lambda_c = 128$ and $\lambda_p = 256$.
\end{proof}

\subsection{Memory-hardness lower bound}\label{sec:security-mh}

\begin{proposition}\label{prop:mh}
Any implementation of \textsc{B3PoW-Scratch} that follows the
specification must perform read and write traffic of magnitude
$\Theta(\texttt{ITERATIONS}\cdot\texttt{LANES}\cdot
\texttt{BLOCK\_BYTES})$ to its scratchpad, or pay a multiplicative
recompute penalty roughly proportional to the size of the omitted
working set.
\end{proposition}

\begin{proof}[Sketch]
The address derivation at iteration $r$ depends on the lane state at
iteration $r$. The lane state at iteration $r{+}1$ depends on the
block read at iteration $r$, and that block depends on the writeback
of every earlier iteration whose derived address coincides with
$\textsc{addr}_r[L]$. An attacker that stores $M < \texttt{SCRATCH\_BYTES}$
bytes of pad and recomputes the missing blocks on demand must, for
every cache miss, redo the work of every prior iteration in that
lane that participated in the missing block's writeback chain. The
fraction of misses is $1 - M / \texttt{SCRATCH\_BYTES}$ in
expectation under address uniformity (\S\ref{sec:security-uniformity}),
giving a multiplicative penalty roughly proportional to
$\texttt{SCRATCH\_BYTES} / M$. This is the same fundamental argument
used by scrypt~\cite{percival2009scrypt} and Argon2d~\cite{biryukov2016argon2},
specialised to a fixed-size on-chip-fitting pad and eight-lane
parallelism. The bound is in the same family as Argon2's but is
weaker in absolute terms (full Argon2 is parameterised on much
larger working sets, at correspondingly worse verification cost).
\end{proof}

\subsection{Verification-DoS resistance}\label{sec:security-verifydos}

Verification of one PoW hash requires the full $\approx 16$\,MiB
memory traffic of the algorithm. We bound the attack surface in
three layers, in part following the CVE-2018-17144 design
discipline~\cite{cve2018-17144} of validating cheap checks first:

\begin{itemize}\itemsep0.2em
\item \emph{Headers-first sync.} The node validates PoW on a header
      before requesting the block body, exactly as in
      Bitcoin Core.
\item \emph{Per-parent pad cache.} The 1\,MiB pad depends only on
      $p$; an LRU keyed on $p$ amortises init across sibling
      headers. The cache stores the pristine init-only pad and
      hands out a fresh copy per call, because the inner loop
      mutates the pad.
\item \emph{Per-hash wall-clock budget.} The consensus
      \texttt{b3pow\_verify\_budget\_ms = 50} bounds the work
      per header; headers that exhaust it are routed to peer
      scoring as \texttt{BLOCK\_POW\_BUDGET}.
\end{itemize}

The combination yields a target p95 verification cost well under 50\,ms
on a 2026-era CPU core, satisfying~G4. We have not yet measured the
final p95; the bench harness in the repository
(\texttt{contrib/testing/bench/}) is the scaffold.

\subsection{Address uniformity}\label{sec:security-uniformity}

The memory-hardness argument (Proposition~\ref{prop:mh}) requires
that the per-lane address sequence is statistically uniform over
$[0, 2\,048)$. We rely on three properties:

\begin{enumerate}\itemsep0.2em
\item the wyhash secret-table multipliers~\cite{wyhash} pass the
      standard PRNG suites for full 64-bit avalanche;
\item the multiply--rotate--XOR pattern in
      Algorithm~\ref{alg:addr} is a standard wyhash-style mixing
      kernel and preserves uniformity from any high-entropy
      lane-state distribution; and
\item a per-lane $\chi^2$ test on $2^{20}$ samples per lane is wired
      into the project's continuous-integration gate
      (\texttt{ref/tests/test\_address\_uniformity.py}).
\end{enumerate}

The v1.1.1 release (\texttt{SPEC\_VERSION = 0x00010101}) restored
distinctness of \texttt{ITER\_MUL}[7] after a duplicate of
\texttt{ITER\_MUL}[1] was discovered in v1.1.0; with the duplicate
in place lane~7's address derivation collapsed to lane~1's and
created a correlated-write pattern in the pad (finding F-1
in~\cite{b3chain-51attack}).

\subsection{Inter-lane diffusion}\label{sec:security-diff}

\begin{lemma}\label{lem:diffusion}
Under the lane permutation $\sigma\colon L\mapsto (5L+1)\bmod 8$
(an 8-cycle in $S_8$), every output lane state in iteration $r$
depends on every input lane state in iteration $r-8$. Combined
with the full intra-lane diffusion of $\texttt{INNER\_ROUNDS}=2$
BLAKE3 rounds on the cv--message pair, the algorithm achieves full
output-on-input diffusion within eight outer iterations, and
maintains it for the remaining 2\,040.
\end{lemma}

\begin{proof}[Sketch]
$\sigma$ decomposes as the single 8-cycle $(0\;1\;6\;7\;4\;5\;2\;3)$,
so $r$ applications of $\sigma$ move every lane by $r$ positions
along the cycle. After eight applications every position has been
visited; combined with BLAKE3's per-lane diffusion bound, this
yields full output-on-input dependence at iteration~8. The argument
generalises to algebraic-degree bounds on $\textsc{MixStep}
\circ \sigma$; a formal proof in the indifferentiability framework
is on the audit scope~\cite{b3chain-auditscope}.
\end{proof}

\subsection{51\,\% attack cost}\label{sec:security-51}

The standard 51\,\%~chain-reorganisation cost~\cite{nakamoto2008bitcoin}
is roughly $H_\text{atk}\cdot t\cdot c_\text{hash}$ where $c_\text{hash}$
is the operating cost of the cheapest unit of hashrate. On
B3PoW-Scratch, that cheapest unit is the reference FPGA card; the
attacker pays the same per-unit cost as the honest side. Selfish-mining
analyses~\cite{eyal2014selfish,bonneau2015sok,sapirshtein2016selfish}
lower the practical threshold below 50\,\% to $\approx 25$--$33\%$
of network hashrate, depending on the propagation advantage
parameter~$\gamma$. The project's 51\,\%-attack
analysis~\cite{b3chain-51attack} catalogues twelve attack vectors,
the calibrated mitigations
(LWMA-3~\cite{zawy2018lwma} difficulty, BIP-94~\cite{bip94} timewarp
enforcement, a consensus \texttt{max\_reorg\_depth} cap, a depth-aware
ban-score in the headers handler, and a two-tier pinned LRU pad
cache), and the per-scenario costs across launch, growth, and
maturity hashrate regimes. We summarise only the conclusion: under
the stated mitigations the dominant residual exposure is the
bootstrap-phase low-hashrate window, and the asymptotic exposure is
the (open-ended) question of an ASIC monopolist.

\subsection{Long-range and time-warp}\label{sec:security-lrtw}

The chain inherits Bitcoin's headers-first sync, sub-second clock
checks, and most-cumulative-work fork-choice rule, so long-range
attacks are bounded by the same arguments as in
Bitcoin~\cite{nakamoto2008bitcoin}. Time-warp is closed on the
mainnet difficulty algorithm by enforcing
BIP-94~\cite{bip94}; the implementation is in
\texttt{src/kernel/chainparams.cpp} of the b3chain repository and
discussed in the 51\,\%-attack analysis as
mitigation~M-2~\cite{b3chain-51attack}.

% --- 6. Hardware analysis -------------------------------------------------
\section{Hardware Analysis}\label{sec:hardware}

The full quantitative treatment is in
\texttt{doc/analysis/FPGA-FEASIBILITY.md} and
\texttt{doc/analysis/ASIC-ECONOMICS.md} in the
repository~\cite{b3chain-fpga,b3chain-asic}; this section gives the
back-of-envelope reasoning that motivates the parameter choices in
Table~\ref{tab:params}.

\paragraph{CPU.}
A reference Python implementation runs at single-digit hashes per
second on a 2026-era Zen~4 core; an optimised C++ AVX-512
implementation is expected in the low-hundreds-H/s range. Neither
is economical at any electricity price.

\paragraph{GPU.}
A modern consumer GPU (e.g.~NVIDIA RTX 4090) has $\sim 128$\,KiB
per-SM L1 cache and a comparable per-warp L2 working set; the
1\,MiB B3PoW-Scratch scratchpad spills to higher cache levels per
nonce. Worse, the sequential RMW chain means nonce $n$ and nonce
$n{+}1$ cannot share work after $\textsc{InitLanes}$; the GPU's only
path to throughput~--- many parallel nonces in flight~--- is exactly
the path the algorithm forecloses. The expected operating point is
tens to low-hundreds of H/s at hundreds of watts, i.e.~below the
CPU H/s\,/\,\$ envelope. The project ships no GPU miner; the
historical kernel~\cite{b3chain-gpuminer} targets a retired
double-BLAKE3 PoW and is retained only as a reference of an
inefficient implementation.

\paragraph{FPGA.}
The reference target is the Xilinx Kintex UltraScale+
XCKU5P~\cite{xilinx-ku5p}. Resource accounting at the algorithm
level is summarised in Table~\ref{tab:fpga}; the dominant resource
is BRAM (one partition per lane), and the LUT and DSP budgets are
modest. At 250\,MHz and one iteration per cycle per pipeline, one
pipeline produces $\approx 1$\,MH/s
($\texttt{LANES} \cdot 250\cdot 10^6 / \texttt{ITERATIONS}$); the
BRAM budget supports eight pipelines, for a target of
$\approx 8$\,MH/s per KU5P card at $\approx 10$\,W typical. Numbers
are pre-measurement.

\begin{table}[t]
\centering
\caption{Reference KU5P resource budget. Source:
         FPGA feasibility note~\cite{b3chain-fpga}, SPEC~\S3
         and~\S8.D~\cite{b3chain-spec}.}
\label{tab:fpga}
\begin{tabular}{lll}
\toprule
Resource          & Available & Used by B3PoW-Scratch \\
\midrule
BRAM (36\,Kb tiles) & 432    & $\approx 150$ (scratchpad partitions) \\
LUTs                & 216k   & $\approx 10$--$30$k (mix pipelines + glue) \\
DSP48E2             & 1\,824 & $8$--$16$ (per-lane $64\times 64$ multiplier) \\
GTH transceivers    & 16     & 0 (Stratum hosted by ESP32-S3) \\
\bottomrule
\end{tabular}
\end{table}

\paragraph{ASIC.}
A custom B3PoW-Scratch ASIC can: pack the scratchpad in dense SRAM
of comparable density to the FPGA's BRAM, eliminate the FPGA's LUT
fabric overhead, run the inner loop at $\sim 1$\,GHz instead of
$\sim 250$\,MHz, and pipeline more lanes in parallel. An optimistic
scaling argument puts a 7\,nm tape-out at roughly $10$--$30\times$
the H/s\,/\,W of the reference FPGA. Crucially, the per-die cost is
dominated by SRAM area, which is comparable across modern fabs;
combined with multi-million-USD NRE, this caps the economical
operating point at network revenues well above the launch hashrate
regime. The full crossover analysis (NRE \emph{vs} network reward
\emph{vs} mass-production die cost) is in the project's ASIC
economics note~\cite{b3chain-asic}; the conclusion is that an ASIC
is possible at sufficient network value, not that it is impossible.

\paragraph{Ranking.}
Table~\ref{tab:rank} summarises the hardware ranking in
order-of-magnitude H/s\,/\,\$ relative to the reference FPGA.
Numbers are pre-measurement; the in-tree benchmark harness will
publish measured values in
\texttt{contrib/testing/bench/results/r0/}.

\begin{table}[t]
\centering
\caption{Order-of-magnitude hardware ranking. Estimates pending
         the round-0 bench results.}
\label{tab:rank}
\begin{tabular}{llll}
\toprule
Tier & Hardware                & H/s/\$ (rel.) & Notes \\
\midrule
1 & Custom 7\,nm ASIC          & $\sim 5$--$30\times$ & Hypothetical;
                                                        bounded by
                                                        SRAM cost \\
2 & B3Miner-1 (KU5P FPGA)      & $1\times$ (ref.)    & Launch miner \\
3 & High-end AVX-512 CPU       & $\sim 10^{-3}$--$10^{-2}\times$ & --- \\
4 & High-end consumer GPU      & $<$ CPU H/s/\$       & sequential
                                                        RMW chain
                                                        defeats SIMD \\
\bottomrule
\end{tabular}
\end{table}

% --- 7. Implementation & parity -------------------------------------------
\section{Implementation and Parity Testing}\label{sec:impl}

The most common silent-failure mode for new PoW launches is a
mismatch between the formal spec and a single production
implementation. We address this with four in-tree implementations
and a single set of consensus vectors that all four are
continuous-integration-gated to reproduce byte-for-byte:

\begin{itemize}\itemsep0.2em
\item \emph{Python reference} (\texttt{contrib/miner/b3miner-rtl/ref/b3pow\_ref.py}),
      deliberately verbose; the executable spec.
\item \emph{C++ consensus} (\texttt{src/crypto/b3pow\_scratch.cpp}),
      the production validator inside the node binary.
\item \emph{TypeScript pool validator}
      (\texttt{contrib/testnet/pool/src/lib/b3pow-scratch.ts}),
      used by the reference Stratum pool.
\item \emph{SystemVerilog RTL}
      (\texttt{contrib/miner/b3miner-rtl/rtl/}), the FPGA pipeline.
\end{itemize}

The single source of truth is
\texttt{src/test/data/b3pow\_consensus\_vectors.json}, generated by
\texttt{ref/gen\_vectors.py}, including: minimum-difficulty headers,
boundary-target tests, non-trivial parent-hash tests (so the pad
init is exercised), and ``cache-pair'' tests (two nonces against
the same parent, to exercise the pad-cache copy semantics). CI
runs all four implementations against the vectors; divergence fails
the gate.

A modest amount of differential fuzzing against the Python
reference, with the C++ port as the production target, is in scope
for the external audit~\cite{b3chain-auditscope}.

% --- 8. Open questions ----------------------------------------------------
\section{Open Questions}\label{sec:openq}

\paragraph{Long-run security budget post-subsidy decay.} The
block-subsidy schedule is unchanged from Bitcoin; the long-run
security model after the subsidy halvings drop below the
fee-market-supported budget is the same open question for B3Chain
as it is for Bitcoin and is treated in the project's economics
notes~\cite{b3chain-securitybudget}.

\paragraph{Algorithm-parameter refresh.} We do not currently plan
parameter refreshes on a fixed cadence; any future change to
\texttt{SCRATCH\_BYTES}, \texttt{ITERATIONS}, or \texttt{LANES} is
a hard fork. A KU15P-targeted v1.2 variant with a 2\,MiB scratchpad
is reserved in the spec but not scheduled.

\paragraph{Stratum~V2 with Noise XX.} Pool-level identity hardening
via Stratum~V2~\cite{stratumv2} with Noise XX transport encryption is
in scope of the external audit. The reference Stratum~V2 pool
implementation is in
\texttt{contrib/testnet/pool/src/sv2/}~\cite{b3chain-spec}.

\paragraph{Independent hardware feasibility.} We would value
independent review from third-party FPGA and ASIC design houses on
the resource budget in Table~\ref{tab:fpga} and on the
break-even analysis in~\cite{b3chain-asic}.

% --- 9. Conclusion --------------------------------------------------------
\section{Conclusion}\label{sec:conclusion}

We described B3PoW-Scratch v1.1, a memory-hard Proof-of-Work built
on BLAKE3 over a 1\,MiB on-chip-memory-bounded scratchpad with eight
parallel lanes and 2\,048 sequential RMW iterations per hash. The
construction is calibrated to make a small on-chip-memory FPGA card
the most economical production miner at the launch-hashrate regime,
\emph{not} to claim permanent ASIC resistance. We argued
preimage and collision resistance inherit from BLAKE3, the
memory-hardness lower bound is in the same family as that of scrypt
and Argon2d (and weaker in absolute terms), the verification cost
is bounded by a per-hash wall-clock budget combined with a
per-parent pad cache, and the catalogued PoW attack surface
(time-warp, selfish mining, verification-DoS, low-hashrate
bootstrap) is mitigated by a combination of inherited Bitcoin Core
discipline and B3Chain-specific consensus rules. A four-implementation
in-tree parity gate removes the most common silent-failure mode for
new PoW launches. The construction is otherwise Bitcoin: the
remainder of the consensus surface, the wire protocol, and the
monetary policy are unchanged.

% --- Bibliography ---------------------------------------------------------
% To match LNCS style at submission time, replace the next line with:
%   \bibliographystyle{splncs04}
\bibliographystyle{unsrt}
\bibliography{references}

\end{document}
